\input{sol_pre.tex}
\input{q_2025-05-24_2.tex}
\begin{document}
\input{sol_head.tex}

\Savol{1}
$11$ ga köpaytiramiz: $55a+5b=495$ $\Rightarrow$ $11a+b=99$. Raqamlar: $a=9$, $b=0$. \ $2a+b=18$.
\Javob{B}

\Savol{2}
$y$ ning har bir hadi $x$ ning mos hadidan $2$ marta katta: $y=2x$. \ $\dfrac{2x+2022x}x=2024$.
\Javob{A}

\Savol{3}
Yaqinlaşiş tezligi: $\frac14+\frac1{12}=\frac13$ (masofa/soat) $\Rightarrow$ $3$ soat.
\Javob{C}

\Savol{4}
$3k=0{,}9\cdot2s$ $\Rightarrow$ $k=0{,}6s$, ya'ni $60\%$.
\Javob{C}

\Savol{5}
$\dfrac1{\sqrt{n+1}+\sqrt n}=\sqrt{n+1}-\sqrt n$ $\Rightarrow$ yiğindi $\sqrt{25}-\sqrt3=5-\sqrt3\approx3{,}27\in[3;4)$.
\Javob{B}

\Savol{6}
$\dfrac{4\left(\sqrt{13}-3\right)}4-\dfrac{9\left(\sqrt{13}+2\right)}9=\sqrt{13}-3-\sqrt{13}-2=-5$.
\Javob{B}

\Savol{7}
$b_1\left(1+q^3\right)=70$, \ $b_1q(1+q)=60$. Bölamiz: $\dfrac{1-q+q^2}q=\dfrac76$ $\Rightarrow$ $6q^2-13q+6=0$ $\Rightarrow$ $q=\frac32$ (ösuvçi).\\
$b_1=\dfrac{60}{\frac32\cdot\frac52}=16$, \ $b_5=16\cdot\dfrac{81}{16}=81$.
\Javob{C}

\Savol{8}
$a_n=124-\frac{10}3(n-1)$ — butun, agar $n-1=3k$: \ $a_n=124-10k>0$ $\Rightarrow$ $k=0,1,\dots,12$ — $13$ ta.
\Javob{B}

\Savol{9}
$\dfrac{2^{16}\cdot5^{10}}{10^8}=\dfrac{2^6\cdot10^{10}}{10^8}=6400$; \ $\dfrac{2^6\cdot5^6}{10^6}=1$. \ $6400-1=6399$.
\Javob{A}

\Savol{10}
$\dfrac{a^4-64}{a^2+8}=a^2-8$, \ $\dfrac{a^6-512}{a^4+8a^2+64}=\dfrac{\left(a^2-8\right)\left(a^4+8a^2+64\right)}{a^4+8a^2+64}=a^2-8$. \ Ayirma $0$.
\Javob{D}

\Savol{11}
$\arcsin\left(-\frac12\right)=-\frac\pi6$, \ $\tg\left(-\frac\pi6\right)=-\frac1{\sqrt3}$.
\Javob{C}

\Savol{12}
$p=\sqrt7+1$, $q=\sqrt5+\sqrt3$: \ ifoda $=\left(\sqrt5-\sqrt3\right)\left(\dfrac1{p-q}-\dfrac1{p+q}\right)=\dfrac{\left(\sqrt5-\sqrt3\right)\cdot2q}{p^2-q^2}$.\\
$\left(\sqrt5-\sqrt3\right)q=2$, \ $p^2-q^2=\left(8+2\sqrt7\right)-\left(8+2\sqrt{15}\right)=2\left(\sqrt7-\sqrt{15}\right)$:
\[\frac4{2\left(\sqrt7-\sqrt{15}\right)}=\frac{2\left(\sqrt7+\sqrt{15}\right)}{7-15}=-\frac{\sqrt{15}+\sqrt7}4.\]
\Javob{C}

\Savol{13}
$s=\sin4x$, $c=\cos4x$: \ $(s+c)(s-c)\left(s^2+sc+c^2\right)=\left(s^2-c^2\right)(1+sc)=-\cos8x\left(1+\frac12\sin8x\right)$
\[=-\cos8x-\frac14\sin16x\ \Rightarrow\ a=-16.\]
\Javob{D}

\Savol{14}
$-15<2x+3<15$ $\Rightarrow$ $-9<x<6$: \ $-8,\dots,5$ — $14$ ta.
\Javob{D}

\Savol{15}
$\left(\sqrt2\right)^6=8<9=\left(\sqrt[3]3\right)^6$ $\Rightarrow$ maxraj manfiy. \ Demak $5x-\lg0{,}1=5x+1>0$ $\Rightarrow$ $x>-0{,}2$.\\
$[-10;20]$ da: $0,1,\dots,20$ — $21$ ta.
\Javob{D}

\Savol{16}
$t=\log_xy$: \ $t+\frac1t=2$ $\Rightarrow$ $t=1$, $y=x$ ($x>0$, $x\ne1$). \ $x^2-x=30$ $\Rightarrow$ $x=6$ ($-5$ mos emas). Bitta yeçim.
\Javob{A}

\Savol{17}
$p=\sqrt[3]{6-x}$, $q=\sqrt[3]{x+3}$: tenglama $p^2+q^2-pq=3$, \ $p^3+q^3=9$.\\
$p^3+q^3=(p+q)\left(p^2-pq+q^2\right)$ $\Rightarrow$ $p+q=3$; \ $(p+q)^2-3pq=3$ $\Rightarrow$ $pq=2$ $\Rightarrow$ $\{p;q\}=\{1;2\}$.\\
$p=1$: $x=5$; \ $p=2$: $x=-2$. \ Yiğindi $3$.
\Javob{D}

\Savol{18}
$\left(25-5^x\right)\left(2^x-1\right)\ge0$ $\Rightarrow$ $0\le x\le2$: \ $0+1+2=3$.
\Javob{B}

\Savol{19}
$x^3-2x^2+2x-1=0$ $\Rightarrow$ $(x-1)\left(x^2-x+1\right)=0$ — bitta haqiqiy ildiz.
\Javob{C}

\Savol{20}
Oğiş burçaklari: $\tg\alpha_1=\sqrt3$ $\Rightarrow$ $60^\circ$; \ $\tg\alpha_2=\frac1{\sqrt3}$ $\Rightarrow$ $30^\circ$. \ $a=60^\circ-30^\circ=30^\circ$.
\Javob{B}

\Savol{21}
$f(x+1)=\dfrac{x^2+x}2$, \ $g(1-x)=\dfrac{x-x^2}2$ $\Rightarrow$ yiğindi $x$. \ $x=x^3$ $\Rightarrow$ $x\in\{-1;0;1\}$, yiğindi $0$.
\Javob{D}

\Savol{22}
$\frac\pi4$ atrofida $\log_{\sin x}\cos x\approx1$, $\ln x<0$, demak modul içi manfiy: $f(x)=\dfrac{\ln\cos x}{\ln\sin x}-\ln x$.\\
$L=\ln\frac{\sqrt2}2=-\frac{\ln2}2$:
\[\left(\frac{\ln\cos x}{\ln\sin x}\right)'_{\frac\pi4}=\frac{-\tg x\cdot\ln\sin x-\ln\cos x\cdot\ctg x}{\ln^2\sin x}=\frac{-2L}{L^2}=\frac4{\ln2}=4\log_2e.\]
$f'\left(\frac\pi4\right)=4\log_2e-\frac4\pi=4\left(\log_2e-\frac1\pi\right)$.
\Javob{A}

\Savol{23}
$F(x)=2\arcsin\dfrac x{\sqrt2}+C$, \ $F(-1)=2\cdot\left(-\frac\pi4\right)+C=0$ $\Rightarrow$ $C=\frac\pi2$.
\Javob{C}

\Savol{24}
Nuqtalar: yuqori $P$, örtadagi $Q$, çapdagi $R$, pastki öng $T$. Böyalgan soha $=\triangle PQT+\triangle RQT$.\\
$\triangle RQT$: asosi $RQ=2$ (gorizontal), balandligi $1$: $S=1$.\\
$\triangle PQT$: uni öz içiga olgan $4\times4$ kvadratdan ($16$) uçta töğri burçakli uçburçak (katetlari $4,4$; \ $1,3$; \ $3,1$) va bitta katakni ayiramiz:
\[S_{PQT}=16-8-1{,}5-1{,}5-1=4.\]
Jami $4+1=5$.
\Javob{A}

\Savol{25}
$CB=AC=1$, \ $BD=CB\ctg30^\circ=\sqrt3$. $AC\parallel BD$ (ikkalasi $CB$ ga perpendikulyar) $\Rightarrow$ $\triangle EAC\sim\triangle EBD$:
\[\frac{AE}{EB}=\frac{AC}{BD}=\frac1{\sqrt3}\ \Rightarrow\ \frac{EB}{AB}=\frac{\sqrt3}{1+\sqrt3}.\]
$S_{ABD}=\frac12\cdot BD\cdot CB=\frac{\sqrt3}2$, \ $S_{BED}=\frac{\sqrt3}2\cdot\frac{\sqrt3}{1+\sqrt3}=\frac{3\left(\sqrt3-1\right)}4$, \ $S_{ABC}=\frac12$:
\[\frac{S_{BED}}{S_{ABC}}=\frac{3\sqrt3-3}2.\]
\Javob{A}

\Savol{26}
$AC=10$, $BD=10$, $CD=6$ $\Rightarrow$ $BC=8$. \ $S_{ABD}=\frac12\cdot4\cdot8=16$; $AE$ — mediana: $S_{ABE}=8$.
\Javob{D}

\Savol{27}
$\tg\dfrac{180^\circ}n=\dfrac{\sqrt3}3$ $\Rightarrow$ $n=6$. \ $S=pr=\dfrac{6\cdot2\sqrt3}2\cdot3=18\sqrt3$.
\Javob{B}

\Savol{28}
Tomon $8$. $\triangle ABD$: $BD^2=64+25-2\cdot8\cdot5\cdot\frac12=49$, $BD=7$.\\
$S_{ABD}=\frac12\cdot5\cdot8\cdot\frac{\sqrt3}2=10\sqrt3$, $p=10$ $\Rightarrow$ $r_1=\sqrt3$.\\
$S_{BDC}=\frac12\cdot3\cdot8\cdot\frac{\sqrt3}2=6\sqrt3$, $p=9$ $\Rightarrow$ $r_2=\frac{2\sqrt3}3$.
\[r_1+r_2=\frac{5\sqrt3}3=\frac5{\sqrt3}.\]
\Javob{C}

\Savol{29}
$B$ — $DC$ ning örtasi, $ODC$ tomoni $a$: $B\left(\frac{3a}4;\,\frac{a\sqrt3}4\right)=\left(6;2\sqrt3\right)$ $\Rightarrow$ $a=8$.\\
$C(8;0)$, $D\left(4;4\sqrt3\right)$, $A\left(2;2\sqrt3\right)$, $F\left(3;3\sqrt3\right)$: \ $\overrightarrow{CF}=\left(-5;\,3\sqrt3\right)$, köpaytma $-15\sqrt3$.
\Javob{A}

\Savol{30}
$HB=AH\tg30^\circ=\frac5{\sqrt3}$, \ $HC=AH\tg45^\circ=5$. $BC\ge HC-HB$ ($B$, $C$ — $H$ dan çiqqan bir nurda bölganda):
\[BC_{\min}=5-\frac5{\sqrt3}=\frac{5\left(\sqrt3-1\right)}{\sqrt3}.\]
\Javob{B}

\Savol{31}
$n(A)=8$, \ $n(B)=5$, \ $n(C)=6$: \ $8+5-6=7$.
\Javob{D}

\Savol{32}
\[P=\frac{C_{22}^2\cdot C_8^1}{C_{30}^3}=\frac{231\cdot8}{4060}=\frac{66}{145}.\]
\Javob{B}

\newpage
\Savol{33}
Asos tomoni $a$. Apofema proyeksiyasi $\frac a2$: balandlik $h=\frac a2\tg60^\circ=\frac{a\sqrt3}2$, apofema $l=\frac{a/2}{\cos60^\circ}=a$.\\
Şar markazi balandlikda, $60^\circ$ li burçak bissektrisasida: $r=\frac a2\tg30^\circ=\frac a{2\sqrt3}$.

\textbf{33.} $\dfrac{2r}h=\dfrac{a/\sqrt3}{a\sqrt3/2}=\dfrac23$ — \textbf{B}.

\textbf{34.} $S_{\text{şar}}=4\pi r^2=\dfrac{\pi a^2}3$, \ $S_{\text{yon}}=4\cdot\frac12\cdot a\cdot a=2a^2$: \ nisbat $\dfrac\pi6$ — \textbf{A}.

\textbf{35.} $V_{\text{şar}}=\dfrac43\pi\cdot\dfrac{a^3}{24\sqrt3}=\dfrac{\pi a^3}{18\sqrt3}$, \ $V_{\text{pir}}=\dfrac13a^2\cdot\dfrac{a\sqrt3}2=\dfrac{a^3\sqrt3}6$: \ nisbat $\dfrac{6\pi}{18\cdot3}=\dfrac\pi9$ — \textbf{F}.
\Javob{33 — B, \ 34 — A, \ 35 — F}

\newpage
\Savol{36}
$x>2$: \ $\lg\dfrac x{x-2}=m$, \ $\dfrac x{x-2}=1+\dfrac2{x-2}$ qiymatlari $(1;\infty)$ ni töliq qamraydi.\\
Demak $m=\lg\frac x{x-2}\in(0;\infty)$ — har biri uçun yagona yeçim, $m\le0$ da yeçim yöq.

a) $m_{\max}=0$. \quad b) $m\in(0;\infty)$.
\Javob{a) 0; \quad b) $(0;\infty)$}

\newpage
\Savol{37}
$\sin^4x+\cos^4x=1-2\sin^2x\cos^2x=1-\frac12\sin^22x\in\left[\frac12;1\right]$.

a) $a\in\left[\frac12;1\right]$ da yeçim bor; eng kiçik natural $a$ (yeçimsiz): $2$.

b) Eng katta butun $a=1$: \ $\sin2x=0$ $\Rightarrow$ $x=\frac{\pi k}2$. \ $\left[-\frac\pi2;\frac\pi2\right]$: $-\frac\pi2,\ 0,\ \frac\pi2$ — $3$ ta.
\Javob{a) 2; \quad b) 3}

\newpage
\Savol{38}
$x\to\frac1x$: \ $f\left(\frac1x\right)+3f(x)=\frac4x$. Ikkinçisini $3$ ga köpaytirib, birinçisini ayiramiz:
\[8f(x)=\frac{12}x-4x\ \Rightarrow\ f(x)=\frac12\left(\frac3x-x\right).\]
a) $f(2)=\frac12\left(\frac32-2\right)=-\frac14$.

b) $\frac3x=x$ $\Rightarrow$ $x=\pm\sqrt3$ — $2$ ta.
\Javob{a) $-\frac14$; \quad b) 2}

\newpage
\Savol{39}
Urinma $(1;2)$ va $(3;0)$ dan ötadi: $k=\frac{0-2}{3-1}=-1$. Demak $f(1)=2$, $f'(1)=-1$.

a) $f'(1)=-1$.

b) $u(x)=x^3-2x^2+x+3$: $u(1)=3$, $u'(x)=3x^2-4x+1$, $u'(1)=0$.
\[h'(1)=u'(1)f(1)+u(1)f'(1)=0\cdot2+3\cdot(-1)=-3.\]
\Javob{a) $-1$; \quad b) $-3$}

\newpage
\Savol{40}
Kesişiş: $\frac{x^2}2=2$ $\Rightarrow$ $x=\pm2$.

a) $S=\displaystyle\int_{-2}^2\left(2-\frac{x^2}2\right)dx=8-\frac{8}3=\frac{16}3$.

b) $Oy$ atrofida: $y\in[0;2]$ da kesim — doira, radiusi kvadrati $x^2=2y$:
\[V=\pi\int_0^22y\,dy=\pi y^2\Big|_0^2=4\pi.\]
\Javob{a) $\frac{16}3$; \quad b) $4\pi$}

\newpage
\Savol{41}
\fig{\begin{tikzpicture}[scale=0.42,font=\small]
\coordinate (O) at (0,0); \coordinate (M) at (-7.071,0); \coordinate (N) at (7.071,0); \coordinate (A) at (-6.325,3.162); \coordinate (B) at (6.325,3.162); \coordinate (C) at (-5.478,0); \coordinate (P) at (0,3.162); \coordinate (H) at (-5.478,3.162);
\draw (N) arc (0:180:7.071); \draw (M)--(N); \draw (A)--(B)--(C)--cycle;
\draw[axs] (O)--(B); \draw[ax] (O)--(P); \draw[ax] (C)--(H);
\fill (O) circle (3pt);
\node[above left] at (A) {$A$}; \node[above right] at (B) {$B$}; \node[below] at (C) {$C$}; \node[below] at (O) {$O$}; \node[below left] at (M) {$M$}; \node[below right] at (N) {$N$}; \node[aux,above] at (P) {$P$}; \node[aux,above] at (H) {$H$};
\end{tikzpicture}}
a) $AB=c$: \ $AC=c\sin15^\circ$, $BC=c\cos15^\circ$:
\[S=\frac12c^2\sin15^\circ\cos15^\circ=\frac{c^2}4\sin30^\circ=\frac{c^2}8=20\ \Rightarrow\ c=4\sqrt{10}.\]
b) $CH\perp AB$: \ $CH=\dfrac{AC\cdot BC}{AB}=c\sin15^\circ\cos15^\circ=\dfrac c4=\sqrt{10}$ — bu $AB$ vatardan $MN$ gaça masofa.\\
$OP\perp AB$ $\Rightarrow$ $PB=\frac c2=2\sqrt{10}$, $OP=\sqrt{10}$:
\[ON=OB=\sqrt{40+10}=5\sqrt2.\]
\Javob{a) $4\sqrt{10}$; \quad b) $5\sqrt2$}

\newpage
\Savol{42}
Kvadrat tomoni $s$. Kataklar böyiça (Pifagor): $BE^2=\left(\frac32s\right)^2+s^2=\frac{13}4s^2$, \ $BM^2=(2s)^2+s^2=5s^2$, \ $EM^2=\left(\frac s2\right)^2+(2s)^2=\frac{17}4s^2$.

a) Kosinuslar teoremasi:
\[\cos a=\frac{\frac{13}4+5-\frac{17}4}{2\cdot\frac{\sqrt{13}}2\cdot\sqrt5}=\frac4{\sqrt{65}}.\]
b) $\sin a=\sqrt{1-\frac{16}{65}}=\frac7{\sqrt{65}}$; \ $s=6$: $BE=3\sqrt{13}$, $BM=6\sqrt5$:
\[S_{BEM}=\frac12\cdot3\sqrt{13}\cdot6\sqrt5\cdot\frac7{\sqrt{65}}=63.\]
\Javob{a) $\frac4{\sqrt{65}}$; \quad b) 63}

\newpage
\Savol{43}
Tomon $s$, $S_{\text{olti}}=\frac{3\sqrt3}2s^2$. $CD$ ga parallel $FA$, ular orasidagi masofa $s\sqrt3$ — bu $A$ dan $CD$ gaça masofa.\\
$B$ dan $CD$ gaça: $\triangle BCD$ da $\angle BCD=120^\circ$ $\Rightarrow$ $BC\sin60^\circ=\frac{s\sqrt3}2$.\\
$G$ — $AB$ ning örtasi: $d(G;CD)=\frac12\left(s\sqrt3+\frac{s\sqrt3}2\right)=\frac{3\sqrt3}4s$, \ $S_{GCD}=\frac12\cdot s\cdot\frac{3\sqrt3}4s=\frac{3\sqrt3}8s^2$.

a) $\dfrac{S_{GCD}}{S_{\text{olti}}-S_{GCD}}=\dfrac{\frac38}{\frac32-\frac38}=\dfrac13$.

b) $s=8$: \ $S_{GCD}=\frac{3\sqrt3}8\cdot64=24\sqrt3$.
\Javob{a) $\frac13$; \quad b) $24\sqrt3$}

\newpage
\Savol{44}
a) Yoyilma — yarim doira: $\pi l=2\pi r$ $\Rightarrow$ $l=2r=6$.

b) $h=\sqrt{36-9}=3\sqrt3$. Silindr radiusi $\rho$, balandligi (öxşaşlikdan) $H=h\left(1-\frac\rho3\right)=\sqrt3(3-\rho)$:
\[S(\rho)=2\pi\rho\cdot\sqrt3(3-\rho),\qquad S'=2\sqrt3\pi(3-2\rho)=0\ \Rightarrow\ \rho=\frac32,\ H=\frac{3\sqrt3}2.\]
$S=2\pi\cdot\frac32\cdot\frac{3\sqrt3}2=\dfrac{9\sqrt3}2\pi$.
\Javob{a) 6; \quad b) $\frac{9\sqrt3}2\pi$}

\newpage
\Savol{45}
a) $S(x)=x(12-2x)$, \ $S'=12-4x=0$ $\Rightarrow$ $x=3$ ($S=18$).

b) Yon tomonlar $y$, asos bilan burçak $\varphi$: kesim — trapetsiya, pastki asos $12-2y$, yuqori asos $12-2y+2y\cos\varphi$, balandlik $y\sin\varphi$:
\[S=(12-2y+y\cos\varphi)\,y\sin\varphi.\]
$\varphi$ tayinlangan: $S'_y=(12-4y+2y\cos\varphi)\sin\varphi=0$ $\Rightarrow$ $y=\dfrac{12}{4-2\cos\varphi}$, u holda $12-2y+y\cos\varphi=6$:
\[S=6y\sin\varphi=\frac{36\sin\varphi}{2-\cos\varphi},\qquad \left(\frac{\sin\varphi}{2-\cos\varphi}\right)'=\frac{2\cos\varphi-1}{(2-\cos\varphi)^2}=0\ \Rightarrow\ \varphi=60^\circ.\]
$y=4$, \ $S_{\max}=\dfrac{36\cdot\frac{\sqrt3}2}{\frac32}=12\sqrt3$.
\Javob{a) 3; \quad b) $12\sqrt3$}

\ShowJavoblar
\end{document}
